\section{Methods}
\label{sec:methods}

\paragraph{Two matched tracks, eight arms.}
Table~\ref{tab:arms-short} lists the arms. Within each track, the
architecture and parameter count stay fixed and only the initialization
changes. Pretrained checkpoints transfer the encoder only, with source
heads discarded. Random floors use the identical target architecture. End
to end, the implemented detectors contain 89{,}996{,}801 (ViT) and
93{,}988{,}865 (CNN) trainable parameters.

\begin{table}[t]
\centering
\caption{Controlled initialization roles. All non-random entries are
released, revision-pinned checkpoints; task heads are discarded.}
\label{tab:arms-short}
\small
\setlength{\tabcolsep}{5pt}
\begin{tabular}{@{}lll@{}}
\toprule
Role & ViT-B/16 track & ConvNeXt-V2-Base track \\
\midrule
Random & fresh initialization & fresh initialization \\
Optical & SatDINO, DINO on fMoW-RGB & BigEarthNet-v2 Sentinel-2 \\
SAR & SARMAE, MAE on SAR-1M & BigEarthNet-v2 Sentinel-1 \\
ImageNet & supervised AugReg & FCMAE $+$ supervised \\
\bottomrule
\end{tabular}
\end{table}

\paragraph{Shared point detector.}
Every arm predicts a stride-4 vessel-presence heatmap from the same
three-channel input (Fig.~\ref{fig:detector}, Appendix~\ref{app:figures}).
Gaussian targets with standard deviation two output pixels mark positives;
LOW-confidence annotations define ignore neighborhoods. Training uses the
penalty-reduced focal objective of CenterNet-style point
detection~\cite{zhou2019objects}. The ViT adapter receives stride-16
features and the CNN adapter stride-32 features with one learned upsampling
stage; both terminate in the same 256-channel head. At inference, local
peaks above a permissive 0.05 candidate floor map to scene coordinates, and
120\,m distance non-maximum suppression removes duplicates. The tiler cuts
whole scenes into 512-pixel windows at stride 384.

\paragraph{Uniform fine-tuning recipe.}
One PyTorch Lightning module, data module, and trainer configuration serve
all cells. AdamW uses learning rate $10^{-4}$, weight decay 0.05,
layer-wise decay 0.65, five warm-up epochs, and cosine decay across at most
50 epochs; early stopping waits four development evaluations. Every cell
uses micro-batch 16, accumulation 1, seed 0, and full-encoder fine-tuning.
The campaign ran on one eight-GPU NVIDIA H100 node; each cell executed as a
single-GPU process in Lightning \texttt{32-true} with CUDA matmul and cuDNN
TF32 disabled. The code itself is hardware-neutral; the uniform hardware
class is a property of the evidence. No arm and no label fraction gets its
own hyperparameters.

\paragraph{Registered contrasts.}
For track $t$, role $r$, and fraction $f$, let $F^{(t)}_{r,f}$ denote
development F1 at that cell's checkpoint-bound operating threshold. The
analysis reports transfer over the track-matched random floor,
\begin{equation}
\Delta^{(t)}_{r,f} = F^{(t)}_{r,f} - F^{(t)}_{\mathrm{random},f},
\label{eq:gain}
\end{equation}
and the within-track SAR--optical contrast,
\begin{equation}
\Delta^{(t)}_{\mathrm{SAR-opt},f} =
F^{(t)}_{\mathrm{SAR},f} - F^{(t)}_{\mathrm{optical},f}.
\label{eq:contrast}
\end{equation}
We assess architecture generality from agreement in the sign and pattern of
Eq.~\eqref{eq:contrast} across tracks, and we never pool the tracks into
one effect estimate. The evidence generator rejects a marker that does not
hash to its cohort binding, a metric inconsistent with its own TP/FP/FN
counts, a best-development value that disagrees with the training curve, a
mixed hardware class, or a partial held-out cohort.
